\section{A. Integral structure, level twelve, and the congruence actions}

We use constant finite group schemes, including in residue characteristic
dividing their order. Thus $C_2$ below means the disjoint union of two
copies of $\operatorname{Spec}\mathbb Z$; the morphism
$C_2\to\mu_2$, $i\mapsto(-1)^i$, will always be written explicitly when
it is used. All scheme exact sequences in this section are exact for the
fppf topology. Our finite-stage notation distinguishes a group from its
schematic envelope by specifying their comparison morphism.

\begin{statement}{A.1}{Proved}{Dictionary and the finite-stage comparison maps}
Fix the following matrices, over $\mathbb Z$:
\[
 \rho=\begin{pmatrix}0&1\\-1&1\end{pmatrix},\qquad
 \sigma=\begin{pmatrix}0&-1\\1&0\end{pmatrix},\qquad
 \tau=\begin{pmatrix}-1&0\\0&1\end{pmatrix},\qquad \omega=-I.
\]
Put
\[
 \lambda_0=\begin{pmatrix}-1&2\\0&-1\end{pmatrix},\quad
 \lambda_1=\begin{pmatrix}-1&0\\-2&-1\end{pmatrix},\quad
 \Pi_0=\langle\lambda_0,\lambda_1\rangle,
 \quad M=\Pi_0^{\mathrm{ab}}=\mathbb Ze_0\oplus\mathbb Ze_1.
\]
The following dictionary fixes the successive versions in \S51.
\end{statement}

\begin{longtable}{@{}p{.18\linewidth}p{.56\linewidth}p{.18\linewidth}@{}}
\toprule
Our symbol & Definition and manuscript symbol & C4 scan; folio\\
\midrule\endhead
$\mathfrak S_{\rm d}$ & $\mathrm{SL}_2(\mathbb Z)/\Pi_0'$, the determinant-one part of the discrete $\mathfrak S$; $\Pi_0'=[\Pi_0,\Pi_0]$ & 33--34; 664--664v\\
$\mathcal S$ & Pushout of $1\to M\to\mathfrak S_{\rm d}\to H\to1$ along $M\to\underline M=\mathbb G_a^2$; his $\underline{\mathcal G}^{+}$ & 34--38; 664--666\\
$H$ & $\mathrm{SL}_2(\mathbb Z)/\Pi_0$, of order $12$; the full determinant-sign quotient has order $24$ & 33, 66; 664, 680v\\
$\mathbb O_N$ & Additive pushout $\mathbb Z_{\rm const}\amalg_{N\mathbb Z}\underline{N\mathbb Z}$, with its ring law in A.5 & 51--54; 673--674v\\
$\gamma$ & $\mathbb O_{12}^{*}$; $\mu:\gamma\to\mathbb G_m$ is evaluation; $\varepsilon_2,\varepsilon_3:\gamma\to C_2$ are the two residue characters & 55--58; 675--676v\\
$M_0$ & $\{c_0e_0+c_1e_1:c_0+c_1\equiv0\pmod3\}$, a lattice with its own vector group $\underline M_0$ & 60--62; 677v--678v\\
$N$ & The corrected $\underline N_{\rho,\sigma}$, with the loop equation in $\underline M_0$ & 62; 678v\\
$Z,G$ & $\underline Z_{\rho,\sigma}$ and $\underline G_{\rho,\sigma}$, namely $\mathbb G_a\rtimes\gamma$ and $\mathcal S\rtimes\gamma$ & 75--76; 685--685v\\
$\Gamma_{\rm e},\mathcal M_{\rm e}$ & The versions of $\Gamma_{\rho,\sigma},\underline M_{\rho,\sigma}$ in (123)--(127) & 70--72; 682v--683v\\
$N_\rho,N_\rho^{?}$ & The smooth closure of the $\rho$-normalizer, and its final bisection (212); A.13 records the smaller literal version (171) & 78--82, 94; 686v--688v, 694v\\
$N_\sigma,N_\sigma^\sharp$ & $\sigma$-normalizer and its intersection with $G^\sharp=\underline M\rtimes\gamma$ & 83--84; 689--689v\\
$\Gamma,\mathcal M$ & The final fibre products (183)--(184), retaining the central $C_2$ & 85--88; 690--691v\\
$\mathcal M(0)$,\newline $\mathcal M(1/2)$,\newline $\mathcal M(-\bar j)$ & The three graph subgroups in A.15 & 88--94; 691v--694v\\
\bottomrule
\end{longtable}

\begin{proof}
The assertions about the discrete groups are proved in A.2. A section
$h\mapsto s_h$ of $\mathfrak S_{\rm d}\to H$ gives the finite-stage
scheme $\coprod_{h\in H}\mathbb A^2_{\mathbb Z}$, with multiplication
\[
 (m,h)(n,k)=(m+a_hn+c(h,k),hk),\qquad
 c(h,k)=s_hs_ks_{hk}^{-1}\in M.
\]
Changing the section to $d_hs_h$ gives the explicit isomorphism
$(m,h)\mapsto(m-d_h,h)$. Thus this construction has a specified,
section-independent meaning. An integral cocycle is obtained from
\[
 \rho\longmapsto(R,0),\quad
 \sigma\longmapsto(S,(1/3,-1/3)),\qquad
 R=\begin{pmatrix}0&-1\\1&-1\end{pmatrix},\quad
 S=\begin{pmatrix}0&1\\1&0\end{pmatrix}.
\]
Choose words $s_h$ lifting the twelve elements of $H$ and let $t_h$ be
their affine translations. Then
$c(h,k)=t_h+a_ht_k-t_{hk}\in\mathbb Z^2$, since the indicated word lies
in $M$. The affine representation kills $\omega$; the retained component
$h$ records it. The pair consisting of the component and affine image is
faithful: on its possible kernel in $M$, the affine image is translation
by the lattice vector itself.

The envelope $E_{\mathfrak S_{\rm d}}$ of \S50 is the inverse limit of
the analogous stages obtained from $nM$. The transition from any of
these stages to the stage above is induced by the inclusion
$nM\hookrightarrow M$ and is a morphism of affine group schemes.
It gives the comparison $E_{\mathfrak S_{\rm d}}\to\mathcal S$; the
formulas below concern this specified finite stage. On profinite
points the continuous map
$\widehat{\mathfrak S_{\rm d}}\to\mathcal S(\widehat{\mathbb Z})$
uses one common $H$-component and the coordinate maps
$\widehat M\to M\otimes\mathbb Z_p$ at every prime. This describes its
image and the meaning of the diagonal points used later in the
manuscript. Likewise $\widehat{\mathbb Z}^{*}\to\gamma(\widehat{\mathbb Z})$
sends $\mu$ to its single residue class $r\bmod12$ and the unique
$k\in\widehat{\mathbb Z}$ with $\mu=r+12k$. The proof of A.5 shows
directly that this is a continuous homomorphism. These comparison maps
retain the integral and profinite data rather than identifying all
points of the schemes with one profinite group.
\end{proof}

\begin{statement}{A.2}{Proved}{The subgroup, its free basis, and the action of the quotient}
The subgroup $\Pi_0$ is normal and free on the displayed
$\lambda_0,\lambda_1$. Its indices in $\mathrm{SL}_2(\mathbb Z)$ and
$\mathrm{GL}_2(\mathbb Z)$ are $12$ and $24$. Moreover
\[
 H\simeq\mathrm{SL}_2(\mathbb Z/4)/V_4\simeq\mathrm{Dic}_3,
 \qquad V_4=\langle\lambda_0\bmod4,\lambda_1\bmod4\rangle.
\]
Conjugation on $M$ sends $\rho,\sigma,\tau$ respectively to $R,S,-I$;
in particular $\tau'=\omega\sigma\tau$ acts as $-S$.
For every divisor $N$ of $24$, the action of
$\mathrm{GL}_2(\mathbb Z)/\Pi_0$ on $M/NM$ has the following orders.
\[
\begin{array}{c|rrrrrrrr}
N&1&2&3&4&6&8&12&24\\\hline
|\operatorname{im}|&1&6&12&12&12&12&12&12\\
|\ker|&24&4&2&2&2&2&2&2
\end{array}
\]
This proves the index and lattice-action portion of Lead 2
(\S51.1--4; C4 scans 34--38, 64--67; fols. 664--666, 679v--681).
\end{statement}
\begin{proof}
Matrix multiplication gives
\[
 \lambda_0=\omega\sigma\rho\sigma\rho,\quad
 \lambda_1=\omega\rho\sigma\rho\sigma,
\]
\[
 \rho\lambda_0\rho^{-1}=\lambda_1,\quad
 \rho\lambda_1\rho^{-1}=(\lambda_1\lambda_0)^{-1},\quad
 \sigma\lambda_0\sigma^{-1}=\lambda_1,\quad
 \sigma\lambda_1\sigma^{-1}=\lambda_0.
\]
Conjugation by $\tau$ inverts each $\lambda_i$. These identities prove
normality and the stated action after abelianization. We use the
amalgam presentation
$\mathrm{SL}_2(\mathbb Z)=C_4*_{C_2}C_6$, with factors generated by
$\sigma,\rho$ and $\sigma^2=\rho^3=\omega$; it is recorded in
\cite[\S4, equation (1)]{A_KirschmerLeedhamGreen}. Adding $\lambda_0=1$
forces $\sigma\rho\sigma^{-1}=\rho^{-1}$; $\lambda_1=1$ is its
conjugate. The resulting group has the twelve normal forms
$\rho^i\sigma^j$, $0\leq i<6$, $0\leq j<2$.
Modulo four, $V_4$ has four elements, is normal, and the images of
$\rho,\sigma$ in $\mathrm{SL}_2(\mathbb Z/4)/V_4$ retain orders six
and four. Since $|\mathrm{SL}_2(\mathbb Z/4)|=48$, this realizes all
twelve normal forms and proves the index claim and the stated quotient.
Surjectivity of the reduction map follows also from elementary row
operations over $\mathbb Z/4$.

The Bass--Serre tree of the amalgam has trivial $\Pi_0$-stabilizers:
each factor, and therefore every conjugate of a factor, injects in
$H$. Its quotient by $\Pi_0$ has $12/4+12/6=5$ vertices and
$12/2=6$ edges. It is connected. Choosing a maximal tree shows directly
that its fundamental group is free of rank $6-5+1=2$, and lifting paths
identifies this group with $\Pi_0$. Since $\lambda_0,\lambda_1$ generate
it, they are a free basis. For completeness, a surjective endomorphism
of a finitely generated free group is injective: free groups are
residually finite by completing the finite labelled path of a nonempty
reduced word to permutations; on the finite set of subgroups of any
fixed finite index, inverse image under a surjection is an injective
permutation. A nontrivial element in its kernel would then lie in every
finite-index subgroup, a contradiction. This applies to the map from
a chosen free basis to the displayed generating pair.

The matrices $R,S$ satisfy $R^3=S^2=(RS)^2=I$ and give the six
distinct matrices of $S_3$ modulo every $N\geq2$. The additional
$-I$ is already $I$ modulo two. For $N\geq3$, $-I$ is outside those
six matrices: direct inspection of
$I,R,R^2,S,RS,R^2S$ proves this over $\mathbb Z/N$.
Thus the image has order twelve for each displayed $N\geq3$.
The central element $\omega$ is the order-two kernel integrally, and
the table follows. The program \texttt{a\_finite.py} constructs the
twenty-four cosets modulo four and propagates the matrix action by a
coupled breadth-first traversal; it also constructs the above quotient
graph and checks its connectivity and Euler characteristic. This is a
finite certificate for the presentation argument, not a finite test
used in place of freeness.
\end{proof}

\begin{statement}{A.3}{Proved}{The precise source of level twelve and the group of order thirty-two}
Let $\Pi_{00}$ be the inverse image in $\Pi_0$ of $M_0$, and put
$\delta=\lambda_1\lambda_0^{-1}$. Then
\[
 \Pi_{00}=\langle\!\langle\delta\rangle\!\rangle_{\mathrm{SL}_2(\mathbb Z)}
 =\Pi_0\cap\mathrm{SL}_2(\mathbb Z)',\qquad
 [\mathrm{SL}_2(\mathbb Z):\Pi_{00}]=36.
\]
Its congruence level is exactly $12$, and the reduction map gives
\[
 K_{32}:=\Pi_{00}/\Gamma(12)
 =V_4\times Q_8
 \ \subset\ \mathrm{SL}_2(\mathbb Z/4)\times\mathrm{SL}_2(\mathbb F_3)
 =\mathrm{SL}_2(\mathbb Z/12).
\]
Here $Q_8$ is the quaternion subgroup of $\mathrm{SL}_2(\mathbb F_3)$.
The subgroup $K_{32}$ is normal in $\mathrm{GL}_2(\mathbb Z/12)$.
This is the group of order $32$ asked for at C4 scan 64, folio 679v;
its exact congruence origin is the additional index-three lattice
$M_0\subset M$, together with the level-four realization of $\Pi_0$.
\end{statement}
\begin{proof}
With $e_i=[\lambda_i]$ we have
\[
 M_0=\mathbb Z(e_1-e_0)+\mathbb Z(3e_1)
     =(1-R)M+(1-S)M.
\]
The normal closure $J$ of $\delta$ is contained in $\Pi_{00}$.
In $\Pi_0/J$, the relation $\lambda_0=\lambda_1$ and its $\rho$-conjugate
give $\lambda_1=(\lambda_1\lambda_0)^{-1}$, so the quotient is cyclic
of order at most three. The map $M\to\mathbb Z/3$, $(c_0,c_1)\mapsto
c_0+c_1$, shows that its order is exactly three. Hence $J=\Pi_{00}$.
Abelianizing the amalgam gives $\mathrm{SL}_2(\mathbb Z)^{\rm ab}
\simeq\mathbb Z/12$, with $[\rho]=2$ and $[\sigma]=3$ after choosing
a generator. Each $[\lambda_i]=4$. The restriction to $\Pi_0$ therefore
has kernel precisely $\Pi_{00}$, proving the intersection formula.
In particular the two occurrences of $12$ in Lead 2 are related by the
exact sequence $0\to C_3\to C_{12}\to H^{\rm ab}=C_4\to0$.

Chinese remainders give the displayed product decomposition. The
matrix of $\delta$ is
\[
 \delta=\begin{pmatrix}1&2\\2&5\end{pmatrix},\qquad
 \delta^2\equiv5I\pmod{12}.
\]
Modulo four its normal closure is $V_4$: conjugation cyclically
permutes the three nonzero elements. Modulo three its normal closure
is the eight-element group consisting of $\pm I$ and the six
trace-zero determinant-one matrices; each of those six squares to
$-I$. It is closed, as one checks on two anticommuting generators,
and is the quaternion group. To see that the normal closure in the
product is the whole product, take commutators of $\delta$ with
$(g,1)$: these generate $V_4\times1$. We can then isolate
$(1,\delta\bmod3)$, whose normal closure is $1\times Q_8$.
Consequently the image of $\Pi_{00}$ is $V_4\times Q_8$.
Its full inverse image has index
$48\cdot24/32=36$, equal to the already proved index of $\Pi_{00}$;
the two subgroups coincide. This proves $\Gamma(12)\subset\Pi_{00}$.

For minimality, the upper transvection
$T^m=\left(\begin{smallmatrix}1&m\\0&1\end{smallmatrix}\right)$
belongs to $\Pi_{00}$ only when $4\mid m$, because the only such
element of $V_4$ is $I$, and only when $3\mid m$, because $Q_8$ has
no nonidentity element of order three. Every $\Gamma(m)$ contains
$T^m$, so $\Gamma(m)\subset\Pi_{00}$ implies $12\mid m$.
Finally $V_4$ is normalized by $\mathrm{GL}_2(\mathbb Z/4)$, and the
trace description of $Q_8$ is invariant under every invertible
conjugator. Their product is normal in the full level-twelve group.
All thirty-two matrices, their orders, their conjugates, and the
normal-closure assertions are exhaustively checked and printed by
\texttt{a\_finite.py}.
\end{proof}

\begin{statement}{A.4}{Proved}{All congruence arrows and the scope of Lead 2}
The exact finite quotients obtained from A.3 have orders
\[
\begin{array}{c|r|r}
\text{group}&\text{order}&\text{quotient by }K_{32}\\\hline
\mathrm{SL}_2(\mathbb Z/12)&1152&36\\
\operatorname{im}(\mathrm{GL}_2(\mathbb Z)\to\mathrm{GL}_2(\mathbb Z/12))&2304&72\\
\mathrm{GL}_2(\mathbb Z/12)&4608&144
\end{array}
\]
The middle image is $\det^{-1}\{1,-1\}$. Thus reduction induces
isomorphisms
\[
 \mathrm{SL}_2(\mathbb Z)/\Pi_{00}\simeq
 \mathrm{SL}_2(\mathbb Z/12)/K_{32},\qquad
 \mathrm{GL}_2(\mathbb Z)/\Pi_{00}\simeq
 \det^{-1}\{1,-1\}/K_{32}.
\]
There is also a surjective homomorphism
$\mathfrak S_{\rm d}\to\mathrm{SL}_2(\mathbb Z/12)/K_{32}$ with kernel
$M_0$. Reduction to $\mathrm{SL}_2(\mathbb Z/12)$ itself does not
factor through $\mathfrak S_{\rm d}$.
These are the actual morphisms relating the action on $M$ to the
order-$32$ subgroup. The action images in A.2 have order at most $12$;
none is the order-$32$ group. They determine the lattice $M_0$ by
$M_0=(1-R)M+(1-S)M$, whose inverse image is used in A.3.
This proves the strongest precise form of Lead 2 and the relation
requested in item 1 (C4 scans 64--67; fols. 679v--681).
\end{statement}
\begin{proof}
The determinant map on $\mathrm{GL}_2(\mathbb Z/12)$ is onto its four
units, with kernel of order $1152$. Integral invertible matrices have
determinant $\pm1$, and the elementary generators of
$\mathrm{SL}_2(\mathbb Z/12)$ lift to integral matrices. Adjoining
$\tau$ therefore gives exactly the middle group. Taking inverse
images of $K_{32}$ proves each isomorphism and each kernel claim.
For the last assertion, direct multiplication gives
\[
 [\lambda_0,\lambda_1]\equiv
 \begin{pmatrix}1&4\\4&5\end{pmatrix}\not\equiv I\pmod{12}.
\]
This element is killed in $\mathfrak S_{\rm d}$ but survives under
reduction, proving the asserted failure of factorization by an
explicit element. It dies after quotient by $K_{32}$, as the exact
kernel calculation requires. This also verifies the qualification in
the margin of scan 67. The scan 66 notation for the Sylow-two group
is transcribed as $D_4$; the six order-four elements and unique
involution prove that its abstract group is $Q_8$, independently of
the notation intended by that handwriting.
\end{proof}

\begin{statement}{A.5}{Proved}{The ring scheme and its exact unit sequence}
For a positive integer $N$, the ring scheme $\mathbb O_N$ is
$\coprod_{r\in\mathbb Z/N}\mathbb A^1_{\mathbb Z}$. A section is locally
written $n+N[x]$, with the equivalence
\[
 (n,x)\sim(n+Nk,x-k)\quad(k\in\mathbb Z).
\]
Addition is the additive pushout law, and multiplication is
\[
 (n+N[x])(m+N[y])=nm+N[ny+mx+Nxy].                    \tag{A.5.1}
\]
Its unit scheme and unit exact sequence are
\[
 \mathbb O_N^*=\coprod_{r\in(\mathbb Z/N)^*}
 \operatorname{Spec}\mathbb Z[t,(r+Nt)^{-1}],
\]
\[
 1\longrightarrow D(1+Nt)\longrightarrow\mathbb O_N^*
 \longrightarrow(\mathbb Z/N)^*_{\rm const}\longrightarrow1. \tag{A.5.2}
\]
The right arrow is smooth and surjective. The corresponding additive
sequence is $0\to\mathbb G_a\to\mathbb O_N\to(\mathbb Z/N)_{\rm const}
\to0$. These prove (64)--(66), (74), and (78)--(79), with the stated
choice of constant finite schemes (C4 scans 51--55; fols. 673--675).

More generally, for a unital associative ring $A$ and a two-sided
ideal $J$ that is finite free as an abelian group, the additive pushout
$A[J]$ has a ring structure and exact sequences
\[
 0\to\underline J\to A[J]\to(A/J)_{\rm const}\to0,
 \quad
 1\to(1+\underline J)^*\to A[J]^*
       \to((A/J)^*)_{\rm const}\to1.                  \tag{A.5.3}
\]
The second surjection is smooth. This supplies the general version of
(69)--(73), including a proof of its local surjectivity.
\end{statement}
\begin{proof}
The pushout is the quotient of
$\mathbb Z_{\rm const}\times\mathbb G_a$ by the constant subgroup
$k\mapsto(Nk,-k)$. Choosing one integer representative of every
residue identifies its components as claimed. Formula (A.5.1) is
unchanged under either replacement $(n,x)\mapsto(n+Nk,x-k)$; direct
expansion proves associativity and distributivity. Equivalently it is
the multiplication inherited from the ring with square-zero-free
ideal coordinate $N\mathbb Z\otimes R$, whose internal product is
$N[x]\,N[y]=N[Nxy]$.

A unit must have unit residue and unit evaluation $n+Nx$.
Conversely choose integers $r,r'$ with $rr'=1+Nh$ and set
\[
 y=-\frac{h+r'x}{r+Nx}.
\]
Then $(r+N[x])(r'+N[y])=1$, proving both sufficiency and the displayed
localized coordinate ring. In residue class one the coordinate is
$1+N[t]$, with multiplication $t\star t'=t+t'+Ntt'$; this is the
kernel in (A.5.2). Every geometric fibre of every component is
nonempty: when the characteristic divides $N$, its representative is
nonzero; otherwise at most one value of $t$ is excluded. The arrow is
an open subscheme of an affine-line bundle over the constant base,
so it is smooth and surjective, and therefore fppf locally admits
sections. This proves exactness without asserting surjectivity on
points over every ring.

For the general construction, take
$(A_{\rm const}\times\underline J)/J_{\rm const}$, where
$j\mapsto(j,-j)$, with multiplication
$(a,x)(b,y)=(ab,ay+xb+xy)$. Each component is an affine space of
rank $\operatorname{rank}J$. Fix a component $b\in(A/J)^*$ and a
lift $a\in A$. On this component multiplication on $J$ on the left
and on the right is a pair of square matrices. The unit locus is
the open set where both determinants are invertible. Indeed, a
chosen lift $a'$ of $b^{-1}$ leaves both inverse errors in $J$;
invertibility of left and right multiplication on $J$ corrects the
two errors, producing a left and a right inverse, which coincide.

To prove that this open set has a point in every geometric fibre,
work over an algebraically closed field $k$. In the finite-dimensional
algebra of pairs of endomorphisms of $J_k$, take the subalgebra
generated by the left-right multiplication pairs from $A$ and $J_k$.
The image of $J_k$ is an ideal, and the image of $a$ is a unit modulo
this ideal. In any finite-dimensional algebra, units lift modulo every
ideal: reduce by the nilpotent Jacobson radical, decompose the
semisimple quotient into matrix blocks, and choose arbitrary invertible
entries in the blocks killed by the ideal; units then lift across the
nilpotent radical by the finite geometric series. Lifting the
adjustment back to $J_k$ gives $a+j$ with both multiplication
operators invertible. The preceding correction supplies its inverse.
Thus every component-unit fibre is nonempty. It is open in affine
space, so the map in (A.5.3) is smooth and surjective. The kernel is
exactly the units of the residue-one component. This also proves the
pair description (73): in the finite free unitization of $J$, a right
inverse is a two-sided inverse, by the determinant criterion for
multiplication on a finite free module.
\end{proof}

\begin{statement}{A.6}{Refuted}{Three literal ring formulas and their exact replacements}
The evaluation criterion in (67b), the multiplication printed in (68),
and the full-kernel assertion (77) are false for the pushout (66).
Their correct replacements are A.5 and
\[
 \ker(\mathbb O_N^*\xrightarrow{\mu}\mathbb G_m)
 =\coprod_{r\in(\mathbb Z/N)^*}
   \operatorname{Spec}\mathbb Z[t]/(r+Nt-1).           \tag{A.6.1}
\]
The residue-one component alone is ${}_N\mathbb G_a$.
After inverting $N$, evaluation and residue give an isomorphism of
ring schemes
$\mathbb O_N\simeq\mathbb G_a\times(\mathbb Z/N)_{\rm const}$,
and of unit schemes
$\mathbb O_N^*\simeq\mathbb G_m\times(\mathbb Z/N)^*_{\rm const}$.
The putative global section of the unit-residue map using the four
integers $\pm1,\pm5$ in (83) is also false over $\mathbb Z$.
These are mathematical corrections to C4 scans 51--56, fols. 673--675v;
they do not settle an unattached handwritten reading.
\end{statement}
\begin{proof}
Over $\mathbb F_2$, the sections $0+2[0]$ and $0+2[1]$ are distinct
points of the same affine component of $\mathbb O_2$, while their
evaluations coincide. Thus evaluation does not detect equality.
For (68), the product $12[1]\,12[1]$ evaluates to $144$ over
$\mathbb Z$; the printed formula evaluates to $12$. Formula (A.5.1)
inserts exactly the missing factor $N$.

Evaluation on the $r$-component is $r+Nt$, giving (A.6.1).
For $N=12$, the section $5+12[-1/3]$ over $\mathbb Q$ has evaluation
one and residue five. It is a nonidentity element of the full kernel,
although ${}_{12}\mathbb G_{a,\mathbb Q}=0$. After $N$ is inverted,
$(\mu,r)\mapsto r+N[(\mu-r)/N]$ is the inverse of the displayed
comparison, proving the stronger structural assertion.
Finally a $\mathbb Z$-point of $\mathbb O_{12}^*$ evaluates to
$\pm1$, and hence has residue $1$ or $11$. The other two residue
components have no integral point. Thus there is no global section
of the four-component residue map over $\mathbb Z$; the numbers
$\pm5$ can be used as affine coordinate representatives only.
All these counterexamples are checked by \texttt{a\_schemes.py}.
\end{proof}

\begin{statement}{A.7}{Proved}{Integral coordinates on the multiplier scheme}
Write $a=\nu_2,b=\nu_3\in\{0,1\}_{\rm const}$, so that
$\varepsilon_2=(-1)^a$ and $\varepsilon_3=(-1)^b$. The four charts of
$\gamma$ have coordinates
\[
 r_0=1-6a+4b,\qquad \mu=r_0+12k\in R^*.
\]
The regular functions used in the manuscript are
\[
 \nu=-3a+2b+6k,\quad
 \nu'=-a+b+3k,\quad \nu''=-a+b+2k,                 \tag{A.7.1}
\]
and satisfy, integrally,
\[
 \mu=1+2\nu=1-2a+4\nu'=1-2b+6\nu'',\qquad
 2\nu'-a=3\nu''-b.                                \tag{A.7.2}
\]
For the alternative representative
$\mu_0=(-5)^a5^b$ used in (84)--(92), with
$\mu=\mu_0+12\nu_0$, the required formulas instead are
\[
 \nu'=-a+b-6ab+3\nu_0,\qquad
 \nu''=-a+b-4ab+2\nu_0.                            \tag{A.7.3}
\]
The transcribed coefficients $-3ab,-2ab$ are false. The substitutions
(A.7.3) correct the corresponding expansions in (138)--(141) as well.
The map $\gamma\to\operatorname{Spec}\mathbb Z$ is smooth, flat, of
relative dimension one, with four geometric components at every prime.
\end{statement}
\begin{proof}
The representatives $1,-5,5,-1$ are congruent to
$(-1)^a\bmod4$ and $(-1)^b\bmod3$. Each equation in (A.7.2) follows
by expansion. Conversely the simultaneous data in (63) give
$k=\nu'-\nu''$, and substitution recovers (A.7.1) without division
by two or three. Thus these are integral isomorphisms of schemes.
Expanding $(-5)^a5^b$ on the four constant components gives
$1-6a+4b-24ab$, whence (A.7.3). The discrepancy is detected over
$\mathbb F_7$ by $a=b=1,\nu_0=0$: $\mu=-25\equiv3$ is a unit,
whereas either of the transcribed formulas reconstructs $\mu=1$.
The exact polynomial checks are in \texttt{a\_schemes.py}.
Smoothness and the component claim follow from the four open
affine-line charts of A.5. In characteristics two and three, each
chart is the entire affine line since its evaluation is a constant
nonzero element. These calculations concern C4 scans 51, 57--58,
74--75; fols. 673, 676--676v, 684v--685.
\end{proof}

\begin{statement}{A.8}{Proved}{The first requested exact sequence, over every base ring}
Set $\delta=e_1-e_0$ and $e=3e_1$, a basis $(\delta,e)$ of $M_0$.
For $\alpha=a_0e_0+a_1e_1$ and $\beta=b_0e_0+b_1e_1$, put
\[
 F(\alpha,\beta)=(1-R)\alpha-(1-S)\beta
   =(-a_0-a_1+b_0-b_1)\delta+a_1e.                 \tag{A.8.1}
\]
Define $N$ by $F(\alpha,\beta)=\nu''e$ \emph{in}
$\underline M_0$, with multiplier in $\gamma$. Then
\[
 (t,s,\underline\mu)\longmapsto
 \left(\underline\mu,
       \alpha=t e_0+\nu''\delta,
       \beta=t e_0+s(e_0+e_1)\right)                \tag{A.8.2}
\]
is an isomorphism of schemes $\mathbb G_a^2\times\gamma\simeq N$,
and its group law is
\[
 (t,s,\underline\mu)(t',s',\underline\mu')
 =(t+\mu t',\ s+\mu s',\ \underline\mu\underline\mu'). \tag{A.8.3}
\]
Consequently the sequence requested at the beginning of \S51.4 is
split exact over $\mathbb Z$:
\[
 1\longrightarrow N^{0+}=\mathbb G_a^2
 \longrightarrow N\longrightarrow\mathbb O_{12}^*
 \longrightarrow1.                               \tag{A.8.4}
\]
The middle group is smooth and flat of relative dimension three,
including at two and three. This proves (104)--(105) after the
author's correction in scans 60--62, fols. 677v--678v.
\end{statement}
\begin{proof}
Multiplication by the displayed $R,S$ gives (A.8.1). The two
coordinates of its required equality say exactly
$a_1=\nu''$ and $a_0+a_1=b_0-b_1$. Therefore the inverse to (A.8.2)
is $t=a_0+a_1$, $s=b_1$, with no division by an integer. In
particular $F:M^2\to M_0$ is a split surjection of free lattices.

We verify the group law and its automorphism meaning. For
$\underline\mu\in\gamma(R)$, define an automorphism $u_{\underline\mu}$
of $\mathcal S_R$ by
\[
 u_{\underline\mu}(\rho)=\lambda_1^{\nu}\rho,\qquad
 u_{\underline\mu}(\sigma)=\omega^a\sigma,
 \qquad u_{\underline\mu}(m)=\mu m\quad(m\in\underline M). \tag{A.8.5}
\]
Exponents on the lattice denote additive vector-group coordinates.
These formulas define the morphism on every component of the
finite-stage scheme. Indeed $(1+R+R^2)e_1=0$ gives
$(\lambda_1^\nu\rho)^3=\omega$, and
$(\omega^a\sigma)^2=\omega$. Substituting in
$\lambda_0=\omega\sigma\rho\sigma\rho$ and
$\lambda_1=\omega\rho\sigma\rho\sigma$ gives multiplication by
$1+2\nu=\mu$ on both lattice generators. The commutator subgroup of
the lattice is therefore killed. These checks construct a compatible
map on the twelve charts of $\mathcal S$, using the scalar map on
their affine coordinates; they do not infer uniqueness merely from
the values of two sections in a nonreduced fibre. The identity
\[
 \nu(\underline\mu\underline\mu')
 =\nu(\underline\mu)+\mu\nu(\underline\mu')
\]
holds as an integral polynomial on every pair of components, by
(A.7.1). The sign $a$ adds modulo two. Thus
$u_{\underline\mu}u_{\underline\mu'}=u_{\underline\mu\underline\mu'}$;
the formula for the inverse follows, so these are automorphisms.

The pair of elements
\[
 \left(\operatorname{Int}(t e_0)u_{\underline\mu},\quad
       (-s(e_0+e_1))u_{\underline\mu}\right)        \tag{A.8.6}
\]
records the effective automorphism and the $\sigma$-normalizer
conjugator of (A.8.2). To check this directly,
$(1-R)\delta=3e_1$ and $\nu=3\nu''-b$ give the first generator
formula after the prescribed $\tau^b$ correction; the second follows
from $t e_0-\beta=-s(e_0+e_1)$ and
$(1-S)(e_0+e_1)=0$. Each factor in (A.8.6) multiplies with scalar
action $\mu$ on its one-dimensional translation coordinate. Their
fibre product over $\gamma$ therefore has exactly (A.8.3), which
proves the transported group law on the manuscript's data. Its
section is $t=s=0$. Smoothness follows from (A.8.2) and A.7.
\end{proof}

\begin{statement}{A.9}{Refuted}{The premature exact sequences and their flat replacement}
The surjection onto $\mathbb G_m\times C_2^2$ in (54), and the
surjection $M^2\to M$ in (59), are false as written. The latter has
cokernel $\mathbb Z/3$, and the correct integral lattice sequence is
\[
 0\to P\to M^2\xrightarrow F M\to\mathbb Z/3\to0,
 \quad P=\mathbb Z(e_0,e_0)\oplus
               \mathbb Z(0,e_0+e_1).              \tag{A.9.1}
\]
The raw scheme obtained by imposing the loop equation only in
$\underline M$ is
\[
 \widetilde Z\simeq
 \operatorname{Spec}_{\gamma}
     \mathcal O_\gamma[t,s,z]/(3z).                \tag{A.9.2}
\]
Its projection to $\gamma$ is split surjective, with kernel
$\mathbb G_a^2\times{}_3\mathbb G_a$ as a scheme. The closed subscheme
$z=0$ is exactly $N$ of A.8 and is the schematic closure of
$\widetilde Z_{\mathbb Z[1/3]}$.
The fine kernel sequence (56), after the sign in (55) is corrected
to $1-R$, is exact with raw kernel
$\mathbb G_a^2\times{}_3\mathbb G_a$ and quotient
${}_4\mathbb G_a$. These are the full corrections to (54)--(59),
(94)--(103) (C4 scans 48--49, 59--62; fols. 671v--672, 677--678v).
\end{statement}
\begin{proof}
The sum of the two coordinates of $F(\alpha,\beta)$ is $3a_1$.
Its image is therefore $M_0$, and (A.8.1) proves surjectivity onto
that lattice and the stated kernel. Tensoring this map with a ring
does not preserve the monomorphism $\underline M_0\to\underline M$:
in its basis $(\delta,e)$ the determinant is $-3$. Imposing the
equation in $\underline M$ instead yields $3(a_1-\nu'')=0$;
the other equation determines $b_0$ from $a_0,a_1,b_1$.
Take $z=a_1-\nu''$, $t=a_0+a_1$, $s=b_1$ to obtain (A.9.2).
Its ideal of functions vanishing after inverting three is $(z)$,
which proves the closure assertion. In particular the characteristic
three raw fibre has dimension four, while the corrected fibre has
dimension three. The kernel version has dimensions three and two.
Finite point enumeration over $\mathbb F_2,\mathbb F_3,\mathbb F_5,
\mathbb F_7$ verifies these counts in \texttt{a\_schemes.py}.

For (54), over the algebraic closure of $\mathbb F_2$, equation (63)
forces $\mu=1$ on every component, so any other nonzero element of
that field has no lift, even fppf locally. Equivalently, the already
specified map to the naive character base factors through $\gamma$;
the correct exact target is the one in A.8. For (56), put
$a=b=0$ and $w=\nu'$. Its equations are $4w=0$ and
$F(\alpha,\beta)=2w e_1$. A section of the projection to
${}_4\mathbb G_a$ is
$\alpha=2w e_0-2w e_1$, $\beta=0$:
$F(\alpha,0)=-6w e_1=2w e_1$ since $4w=0$.
The kernel is $F=0$, already computed. The literal plus sign in
$(1+R)\alpha$ in (55) is inconsistent with (57), for example at
$\alpha=\beta=e_0,w=0$ over $\mathbb Z$. We have stated which
replacement makes the sequence exact.

Scan 61 (100)--(101) contains another sign discrepancy visible in
the attached image: the zero equation (96) requires
$F(\alpha,\beta)-\nu''e=0$, or equivalently
$(1-R)(\alpha-\nu''\delta)-(1-S)\beta=0$.
The plus sign in the shifted $\alpha$ gives $6e_1$ on
$\alpha=\delta,\beta=0,\nu''=1$ over $\mathbb Z$.
Thus this replacement is forced by (94)--(96); it is not an
unannounced choice between readings.
\end{proof}

\begin{statement}{A.10}{Proved}{The integral comparison with the fine scheme of LM1}
The fine scheme of scan 41, with component bits $a,b$, coordinate
$w=\nu'$, and loop equation $3a_1=2w-a+b$, is isomorphic over
$\mathbb Z$ to the corrected $N$ of A.8. The comparison and its
inverse are
\[
 k=w-a_1,\quad \nu''=a_1,\quad
 t=a_0+a_1,\quad s=b_1,
 \qquad
 w=3k-a+b,\quad a_1=2k-a+b,                       \tag{A.10.1}
\]
with $\alpha,\beta$ recovered by (A.8.2). In particular it is an
isomorphism on fibres at two and three and on points over every ring.
The alternate LM1 coordinate using $\nu=2w+a$ is carried to the
present coordinate using $\nu=2w-a$ by $w_{\rm here}=w_{\rm LM1}+a$.
This settles the finite-stage comparison left in LM1 E.20 for this
specific parameter scheme (C4 scans 40--41, 62, 71--72;
fols. 667v--668, 678v, 683--683v).
\end{statement}
\begin{proof}
The equation $2w-3a_1=a-b$ implies
$w=3(w-a_1)-a+b$ and $a_1=2(w-a_1)-a+b$.
Thus (A.10.1) is integral in both directions. Substituting gives
$\mu=1-2a+4w=1-6a+4b+12k$, the same inverted element in both
coordinate rings. The remaining loop coordinate equation is exactly
$a_0+a_1=b_0-b_1$, whose solutions are (A.8.2).
This proves the claimed scheme isomorphism, including all nilpotent
test rings. Both groups use the automorphism and conjugator
multiplication (A.8.6), so the comparison also respects their group
structures. The final coordinate translation follows by equating
$2w_{\rm here}-a=2w_{\rm LM1}+a$ as integral functions; substitution,
not cancellation in a characteristic-two ring, proves the equality.
\end{proof}

\begin{statement}{A.11}{Proved}{Effective automorphisms and the early quotient sequences}
The effective group, its geometric extension, and the early quotient
groups are
\[
 Z=\mathbb G_a\rtimes_\mu\gamma,
 \quad G=\mathcal S\rtimes_u\gamma,
 \quad \Gamma_{\rm e}=\mathbb G_a\rtimes_\mu\gamma,
\]
\[
 \mathcal M_{\rm e}\simeq(\mathbb G_a\times\mathcal S)\rtimes\gamma,
 \qquad
 \mathcal M_{\rm e}^{\sharp}\simeq\mathbb G_a^3\rtimes_\mu\gamma.
                                                               \tag{A.11.1}
\]
The $\gamma$-action on $\mathcal S$ is (A.8.5). In $N$, the subgroup
$L_0$ is the $t$-axis and the quotient is the $s$-coordinate. We obtain
the exact sequences
\[
\begin{gathered}
1\to L_0\to N\to\Gamma_{\rm e}\to1,\qquad
1\to\mathbb G_a\to\Gamma_{\rm e}\to\gamma\to1,\\
1\to\mathcal S\to\mathcal M_{\rm e}\to\Gamma_{\rm e}\to1,
\quad
1\to\underline M\to\mathcal M_{\rm e}^{\sharp}\to\Gamma_{\rm e}\to1,\\
1\to L_0\to Z\to\gamma\to1,
\qquad1\to\mathcal S\to G\to\gamma\to1.
\end{gathered}                                                     \tag{A.11.2}
\]
All are split and exact on points over every ring. The lattice kernel
sequences (127)--(132) are
$0\to\mathbb Z\to P\to\mathbb Z\to0$ and their vector-group
base changes; the maps are $(t,s)\mapsto s$ and $t\mapsto(t,0)$.
These prove (123)--(127), (131)--(132), (142), (145)--(150)
(C4 scans 70--76; fols. 682v--685v).
\end{statement}
\begin{proof}
Under (A.8.6), the first factor is
$\operatorname{Int}(t e_0)u_{\underline\mu}$ and the second has
coordinate $-s$. The kernel of the effective projection is the
$s$-axis; its image is the displayed $Z$. Inner conjugation by
$t e_0$ embeds $L_0\simeq\mathbb G_a$ into $Z$, and (A.8.5)
normalizes it by the scalar $\mu$. Thus the quotients are represented
by the displayed affine schemes, with the stated sections.

For the crossed quotient defining $\mathcal M_{\rm e}$, take
$\mathcal S\rtimes N$ with the action of $N$ through $Z$ and quotient
by $t\mapsto(-t e_0,(t,0,1))$. This is a normal subgroup because
its action on $\mathcal S$ equals inner conjugation by $t e_0$.
Every coset has exactly one representative with $t=0$; its remaining
coordinates are $(x,s,\underline\mu)$ and multiply as in (A.11.1).
The same calculation with $\underline M$ in place of $\mathcal S$
gives the sharp subgroup. All kernel and section assertions now
follow by coordinate projection, so they are exact on arbitrary
ring points and hence as fppf sheaves. The kernel lattice $P$ was
computed in A.9; its quotient by the diagonal $e_0$-line has the
free generator $(0,e_0+e_1)$. This proves every lattice and
vector-group sequence asserted here.

The ambiguous exponent in scan 76 (151) is fixed by the verified
formula $\lambda_1^{\nu}\rho$ in scan 93, and independently by
$u(\lambda_i)=\mu\lambda_i$ in A.8. The candidate exponent
$6\nu''$ fails this test, since the required relation is
$\nu=3\nu''-b$. We retain this correction explicitly in the
readings ledger.
\end{proof}

\begin{statement}{A.12}{Proved}{The two normalizers and the central extension of the multiplier scheme}
Let $d=\nu''$ and set, in $G=\mathcal S\rtimes\gamma$,
\[
 v_{\underline\mu}=\sigma^b\delta^{-d}u_{\underline\mu}.
\]
For the raw normalizer of (152), before the closure in (159),
the literal sequence (155) is
\[
 1\to C_6\times{}_3\mathbb G_a\to N_\rho^{\rm raw}
   \to\gamma\to1,
 \qquad N_\rho^{\rm raw}\simeq{}_3\mathbb G_a\times N_\rho.
                                                               \tag{A.12.0}
\]
The smooth closure of the $\rho$-normalizer is represented by
$\{\rho^k v_{\underline\mu}:k\in\mathbb Z/6\}$; hence it has the
exact sequence
\[
 1\to C_6=L_\rho\to N_\rho\to\gamma\to1.          \tag{A.12.1}
\]
The final subgroup specified in (212), scans 93--94, is
\[
 N_\rho^?=\{\omega^i v_{\underline\mu}:i\in\mathbb Z/2\},
 \qquad 1\to C_2=\mu_\rho\to N_\rho^?\to\gamma\to1. \tag{A.12.2}
\]
Its multiplication on $C_2\times\gamma$ is
\[
 (i,\underline\mu)(i',\underline\mu')
  =(i+i'+c(\underline\mu,\underline\mu'),
       \underline\mu\underline\mu'),\qquad
 c=(a+b)b'\pmod2.                                 \tag{A.12.3}
\]
Thus (A.12.2) is the pullback along
$(\varepsilon_2,\varepsilon_3):\gamma\to C_2^2$ of the order-eight
dihedral central extension $1\to C_2\to D_8\to C_2^2\to1$.
Moreover
\[
 N_\sigma^\sharp=\mathbb G_a\rtimes_\mu\gamma,\qquad
 N_\sigma=(C_4\times\mathbb G_a)\rtimes\gamma,       \tag{A.12.4}
\]
where $\gamma$ acts by $\varepsilon_2$ on $C_4$ and by $\mu$ on
$\mathbb G_a$. These assertions prove (155), (159), and (173)--(181),
and compute the extension asked for after (196)
(C4 scans 78--84, 87, 93--94; fols. 686v--689v, 691, 694--694v).
\end{statement}
\begin{proof}
The identities $(1-R)\delta=3e_1$ and $\nu-3d=-b$ give
\[
 \delta^{-d}u_{\underline\mu}(\rho)\delta^d
   =\lambda_1^{-b}\rho.
\]
For $b=1$, direct multiplication in the modular group gives
$\sigma\lambda_1^{-1}\rho\sigma^{-1}=\rho^{-1}$.
For $b=0$ the displayed expression is $\rho$.
Thus $v_{\underline\mu}$ has the required normalizing action
$\rho\mapsto\rho^{\varepsilon_3}$.
In the finite quotient $H$, the centralizer of $\rho$ is its cyclic
subgroup of order six. Within each of these six components, the
remaining translation condition is $(1-R)m=0$. Its kernel has
equations equivalent to $3z=0$, because $\det(1-R)=3$.
Explicitly $m=z\delta$, with $3z=0$. The cyclic component
$C_6$ fixes this vector. Conjugation by $v_{\underline\mu}$
acts on it as $\mu\varepsilon_3$, and
$\mu\varepsilon_3=1+6\varepsilon_3d$; hence it also fixes
$z\delta$. This proves the central direct product (A.12.0)
and the original sequence (155), including its extra fibre
in characteristic three.
The closure of its characteristic-not-three part is the zero
translation. Translation by $v_{\underline\mu}$ therefore gives
exactly the smooth normalizer stated above, including all its
components, with kernel $L_\rho$.

The identity
\[
 d(\underline\mu\underline\mu')
   =\varepsilon_3(\underline\mu')d(\underline\mu)
      +\mu d(\underline\mu')                     \tag{A.12.5}
\]
is an integral polynomial identity on the sixteen pairs of charts;
it follows by multiplying $\mu=\varepsilon_3+6d$ and using (A.7.1).
In particular its proof does not cancel six in a test ring. Moving
$\sigma^{b'}$ past $\delta^{-d}$ uses $S\delta=-\delta$, and
$u_{\underline\mu}(\sigma)=\omega^a\sigma$. The two central factors
are $\omega^{ab'}$ and $\omega^{bb'}$. Equation (A.12.5) then gives
\[
 v_{\underline\mu}v_{\underline\mu'}
     =\omega^{(a+b)b'}v_{\underline\mu\underline\mu'}.
\]
Consequently the multiplication in $N_\rho$ is
$(k,\underline\mu)(k',\underline\mu')
=(k+\varepsilon_3k'+3c,\underline\mu\underline\mu')$.
Restricting $k$ to $0,3$ proves (A.12.2)--(A.12.3).
The bilinear function $((a,b),(a',b'))\mapsto(a+b)b'$ is a
two-cocycle. The resulting eight-element group has five involutions
and two elements of order four, and is generated by two involutions
with noncommuting product of order four. This proves it is $D_8$;
the full cocycle and element-order checks are in
\texttt{a\_finite.py}. Exactness of (A.12.2) holds even on ring
points, because $v_{\underline\mu}$ is a global section as a scheme,
although it is not a group section.

For $\sigma$, the centralizer in $H$ is $C_4$, generated by
$\sigma$. Its translation kernel is
$\ker(1-S)=\mathbb Z(e_0+e_1)$, also after every base change.
These translations commute with $\sigma$, so the geometric
centralizer is $C_4\times\mathbb G_a$. Equation (A.8.5) supplies
a group section over $\gamma$ and its action on both factors,
which proves (A.12.4) and all the stated exact sequences. The sharp
intersection retains just the translation component. Similarly
$N_\rho\cap G^\sharp$ is the section over
$\gamma'=\ker\varepsilon_3$, since the component
$\rho^k\sigma^b$ is the identity precisely when $k=b=0$.
Put $\gamma''=\ker\varepsilon_2$, $\gamma^0=\gamma'\cap\gamma''$,
and $\gamma'''=\{a=b\}$. The neutral partial section of (166)--(167)
is $v(\gamma^0)$. The partial section of (168)--(169) is exactly
\[
 \underline\mu\longmapsto\omega^b v_{\underline\mu}
 \quad(\underline\mu\in\gamma'''),
\]
because its cocycle vanishes on $a=b$, and at $\underline\mu=-1$
it is $\tau'$. Conjugation by $\tau'$ sends
$v_{\underline\mu}$ on $\gamma'$ to
$\omega^a v_{\underline\mu}$, as (A.12.3) shows. It preserves
the partial section precisely on $\gamma^0$.
This proves the prime and double-prime intersections and the
normalization assertions of (160)--(170) with their exact domains.
\end{proof}

\begin{statement}{A.13}{Refuted}{The literal definition (171) has only half the stated image}
The group defined in (171) from $Z_\rho^{0\sharp}$, $\omega$, and
$\tau'$ maps onto
\[
 \gamma'''=\{\varepsilon_2=\varepsilon_3\}\subset\gamma,
\]
and is the inverse image of this subgroup in the final $N_\rho^?$
of (212). Thus (172) with target all of $\gamma$ is false for
the literal definition (171). Its exact replacement is
$1\to C_2\to N_{\rho,\mathrm{literal}}^?\to\gamma'''\to1$;
the final fibres (212) yield the larger exact sequence (A.12.2).
The surviving last line of (194), with $N_\rho$ and no question mark,
requires the replacement
$1\to C_3\to N_\rho/\mu_\rho\to\gamma\to1$;
the isomorphism to $\gamma$ holds for $N_\rho^?/\mu_\rho$.
The final $\Gamma$ in (184) also satisfies
$\Gamma/C_2\simeq\mathbb G_a\rtimes\gamma$; the unquotiented
version of this isomorphism in the third row of (195) is false.
These corrections concern C4 scans 82, 85--87, 94;
fols. 688v, 690--691, 694v.
\end{statement}
\begin{proof}
The identity component $Z_\rho^{0\sharp}$ maps isomorphically to
$\gamma^0=\{a=b=0\}$ by A.12. The central $\omega$ has trivial
character, and $\tau'$ has multiplier $-1$, hence $(a,b)=(1,1)$.
The literal generated subgroup therefore has image in $a=b$.
Conversely multiplication by $-1$ carries the whole $00$ chart to
the whole $11$ chart, so the image is exactly $\gamma'''$.
At $\underline\mu=-1$, $d=0$ and
$\tau'=\omega v_{-1}$; hence both lifts over each of these charts
occur. This proves the asserted inclusion and its exact image and
kernel. Over $\mathbb F_2$, the $\gamma$-point with $a=1,b=0,k=0$
has evaluation one and has no lift to this literal subgroup. This
finite counterexample is recorded in \texttt{a\_schemes.py}.

Quotienting (A.12.1) by its central order-two subgroup leaves
kernel $C_6/C_2=C_3$, which proves the correction to (194).
The final fibre product defining $\Gamma$ has kernel
$C_2\times\mathbb G_a$ over $\gamma$, by A.12. Its quotient by
the indicated central $C_2$ is $N_\sigma^\sharp$, giving the claimed
isomorphism. On $\mathbb F_2$ points the two groups have $32$ and
$16$ elements, respectively; the explicit projection has kernel of
order two. Thus the missing quotient in (195) is detected by an
actual map and a finite counterexample.
\end{proof}

\begin{statement}{A.14}{Proved}{The final fibre products and every surviving extension sequence}
Using the final (212), define
\[
 \Gamma=N_\rho^?\times_\gamma N_\sigma^\sharp,
 \quad \mathcal M=\Gamma\times_\gamma G,
 \quad \mathcal M^\sharp=\Gamma\times_\gamma G^\sharp.
                                                               \tag{A.14.1}
\]
The coordinates on $\Gamma$ are $(i,s,\underline\mu)$ and its law is
\[
 (i,s,\underline\mu)(i',s',\underline\mu')
  =(i+i'+(a+b)b',\ s+\mu s',\
       \underline\mu\underline\mu').              \tag{A.14.2}
\]
There are central quotient morphisms, with the indicated kernels,
\[
 \Gamma\twoheadrightarrow\Gamma_{\rm e},\qquad
 \mathcal M\twoheadrightarrow\mathcal M_{\rm e},\qquad
 \mathcal M^\sharp\twoheadrightarrow\mathcal M_{\rm e}^\sharp,
 \qquad\ker=C_2=\mu_\rho.                         \tag{A.14.3}
\]
In addition to A.8, A.11, and A.12, the final exact sequences are
\[
\begin{gathered}
1\to\mathcal S\to\mathcal M\to\Gamma\to1,\qquad
1\to\underline M\to\mathcal M^\sharp\to\Gamma\to1,\\
1\to C_2\times\mathbb G_a\times\mathcal S
 \to\mathcal M\to\gamma\to1,\\
1\to C_2\times\mathbb G_a\times\underline M
 \to\mathcal M^\sharp\to\gamma\to1,\qquad
1\to C_2\times\mathbb G_a\to\Gamma\to\gamma\to1.
\end{gathered}                                                     \tag{A.14.4}
\]
For $X=\Gamma,\mathcal M,\mathcal M^\sharp$, the auxiliary group
$S_0X=Z\times_\gamma X$ has split exact sequence
$1\to L_0\to S_0X\to X\to1$.
These prove (182)--(193) and (196), with the corrected quotient
versions of (194)--(195). The central extensions are precisely the
pullbacks of the explicit $D_8$ extension in A.12.
\end{statement}
\begin{proof}
Every map to $\gamma$ in (A.14.1) is a homomorphism, and the
multiplication of the fibre product is componentwise. Insert
(A.12.3)--(A.12.4) to obtain (A.14.2). Forgetting $i$ is a
homomorphism, since its cocycle is central, and its kernel is the
constant $C_2$. The resulting three groups have exactly the
coordinates and actions in (A.11.1), proving the actual comparisons
(A.14.3) between the early and final revisions.
The projection $\mathcal M\to\Gamma$ has the group section
$(a,b)\mapsto(a,b,u_{\underline\mu})$; its kernel is $\mathcal S$.
Replacing $\mathcal S$ by $\underline M$ gives the sharp row.
Kernels over $\gamma$ are the products of the three displayed
geometric kernels, proving (A.14.4). These projections are onto
on ring points: use $v_{\underline\mu}$, $u_{\underline\mu}$, and
zero translation as a scheme section. Finally $Z\to\gamma$ has
the group section (A.8.5), so the auxiliary projections are split
exact. This proves exactness and representability directly for the
entire final diagram.

The central factor $\mu_\rho$ here is $(\omega,1,1)$ in
$(a,b,U)$-coordinates. The central subgroup $\mu_G$ is
$(1,1,\omega)$ and has a separate inclusion. The quotient map
in (A.14.3) removes the first of these, exactly as its kernel states;
no central factors have been combined without a map.
\end{proof}

\begin{statement}{A.15}{Proved}{The three distinguished subgroups: all requested points and fibres}
The final subgroups of scans 88--94 are the graph subgroups
\[
 \mathcal M(0)=\{U=u_{\underline\mu}\},\qquad
 \mathcal M(1/2)=\{U=b_\sigma\},\qquad
 \mathcal M(-\bar j)=\{U=a_\rho\}                 \tag{A.15.1}
\]
inside $\mathcal M=\{(a_\rho,b_\sigma,U)\}$.
Projection on $(a_\rho,b_\sigma)$ identifies each with $\Gamma$ as
a group scheme over $\mathbb Z$. They are smooth and flat of relative
dimension two, with eight geometric components at every prime.

For any connected ring $R$, their common point set is exactly
\[
 \left\{(a,b,k,i,s): a,b,i\in\{0,1\},\ k,s\in R,
           \ 1-6a+4b+12k\in R^*\right\}.          \tag{A.15.2}
\]
For a general ring the bits are locally constant functions on its
spectrum. The group law is (A.14.2), with multiplier product
\[
\begin{split}
 a''&=a+a'\pmod2,\qquad b''=b+b'\pmod2,\\
 k''&=\frac{r_0r_0'-r_0''}{12}
             +r_0k'+r_0'k+12kk'.
\end{split}                                                       \tag{A.15.3}
\]
The numerator is an integer divisible by twelve on every component.
In particular the requested rings have the following complete
parametrizations, with $s\in R$ and $i\in C_2$ always free.
\begin{center}\small
\begin{tabular}{p{.19\linewidth}p{.71\linewidth}}
\toprule
$R$ & Allowed multipliers and component bits; $k=(\mu-r_0)/12$\\\midrule
$\mathbb Z$ & $\mu=1:(a,b,k)=(0,0,0)$; $\mu=-1:(a,b,k)=(1,1,0)$.\\
$\mathbb Z[1/2]$ & $\mu=\pm2^m$ ($m\in\mathbb Z$); $a$ arbitrary,
                   $(-1)^b=\mu\bmod3$.\\
$\mathbb Z[1/3]$ & $\mu=\pm3^m$ ($m\in\mathbb Z$); $b$ arbitrary,
                   $(-1)^a=\mu\bmod4$.\\
$\mathbb Z[1/6]$ & $\mu=\pm2^m3^n$ ($m,n\in\mathbb Z$); $a,b$ arbitrary.\\
$\mathbb F_2$ & $a,b,k$ arbitrary; $\mu=1$;
              $|\mathcal M(p)(\mathbb F_2)|=32$.\\
$\mathbb F_3$ & $a,b,k$ arbitrary; $\mu=(-1)^b$;
              $|\mathcal M(p)(\mathbb F_3)|=72$.\\
\bottomrule
\end{tabular}
\end{center}
Here $p$ denotes any of the three displayed points, not a prime.
As abstract groups, the integral points are
$C_2\times(\mathbb Z\rtimes\{\pm1\})$; the two finite groups are
$D_8\times C_2^2$ over $\mathbb F_2$ and
$C_3\times(C_3\rtimes D_8)$ over $\mathbb F_3$, where $D_8$ acts
on the second $C_3$ through $\varepsilon_3$.
\end{statement}
\begin{proof}
The maps $\Gamma\to G$ given respectively by
$(a_\rho,b_\sigma)\mapsto u_{\underline\mu},b_\sigma,a_\rho$
are homomorphisms. Their graphs have the stated inverse projection,
proving all three group-scheme isomorphisms, not just equality of
dimensions or of point counts. Formula (A.15.2) follows from
$\Gamma\simeq C_2\times\mathbb G_a\times\gamma$ as a scheme.
Equations (A.14.2) and multiplication of
$r_0+12k$ give (A.15.3); its carry term is integral because residue
classes multiply according to the two sign bits.

For $\mathbb Z$, invertibility forces $\mu=\pm1$ and the congruences
in (A.7.2) force exactly the displayed bits. The cocycle then vanishes
since $a=b$. For $\mathbb Z[1/2]$, the unit group is
$\{\pm2^m\}$, and divisibility by twelve in that ring is equivalent
to the displayed congruence modulo three; modulo four imposes no
condition because two is inverted. The argument with the primes
reversed proves the next row, and inverting six proves the fourth.
Uniqueness of $k$ in these torsion-free rings follows directly.
In the two finite fields, twelve vanishes and all representatives
$r_0$ are units, giving $4\cdot2\cdot2\cdot2=32$ and
$4\cdot3\cdot2\cdot3=72$ points. These were also counted directly
from the original lattice equations by \texttt{a\_finite.py}.

There is an explicit ring isomorphism
$\mathbb O_N(\mathbb Z/q)\to\mathbb Z/Nq$,
$r+N[t]\mapsto r+Nt$. Therefore
$\gamma(\mathbb F_2)=(\mathbb Z/24)^*\simeq C_2^3$ and
$\gamma(\mathbb F_3)=(\mathbb Z/36)^*\simeq C_3\times C_2^2$.
The residue maps split as maps of these finite abelian groups, with
kernels $C_2$ and $C_3$. Pulling back the $D_8$ extension yields the
claimed finite-group descriptions, including its action on $s$.
Finally each of the eight charts is an open subset of affine
two-space, and over characteristics two and three it is the whole
affine two-space. This proves smoothness, flatness, and the component
count on the actual fibres.

Under (A.14.3), the first two graphs contain $\mu_\rho$, while the
third intersects it trivially. Indeed $U=a_\rho$ forces
$U=\omega$ when $a_\rho=\omega$. Its image in
$\mathcal M/\mu_\rho$ is a bisection over $\Gamma_{\rm e}$, with
kernel $\mu_G$. This gives explicitly the change of descriptions in
scans 91--94 rather than identifying the central quotients with
all three unquotiented graphs.
\end{proof}

\begin{statement}{A.16}{Proved}{An actual comparison with the groups of \S48}
Let $H_{48}$ be the group of pairs of quadratic-normalizer matrices
\[
 A=c\rho+dI,\quad
 B=\begin{pmatrix}u&t\\-et&eu\end{pmatrix},\quad
 Q_3=c^2+cd+d^2,\quad Q_2=t^2+u^2,
\]
with $Q_2,Q_3$ units and $e^2=(e')^2=1$, satisfying
$\det B=eQ_2=e'Q_3$. Its character map is
\[
 H_{48}\to\gamma_{48}=\mathbb G_m\times\mu_2^2,
 \quad(A,B)\mapsto(\det B,e,e').
\]
There is a specified comparison of character schemes
\[
 \iota:\gamma\longrightarrow\gamma_{48},\qquad
 \underline\mu\longmapsto(\mu,(-1)^a,(-1)^b),      \tag{A.16.1}
\]
which is an isomorphism over $\mathbb Z[1/6]$. Define the actual
base change $H_{48}^\dagger=H_{48}\times_{\gamma_{48}}\gamma$.
Projection followed by the zero-translation section is a morphism
of affine group schemes
\[
 H_{48}^\dagger\longrightarrow\Gamma/\mu_\rho,
 \quad(A,B,\underline\mu)\longmapsto(0,\underline\mu). \tag{A.16.2}
\]
Its kernel is $ST_\rho\times ST_\sigma$, the product of the two
norm-one quadratic groups, and its image is exactly the displayed
copy of $\gamma$. For the full groups the analogous map
$\mathcal M_{48}\times_{\gamma_{48}}\gamma\to\mathcal M/\mu_\rho$
has kernel $ST_\rho\times ST_\sigma\times\mathrm{SL}_2$ and image
the canonical $\gamma$ section.

The \S48 identifications
$\mathcal M_{48}(0)\simeq H_{48}$ and
$\mathcal M_{48}(1/2)\simeq H_{48}$ hold by their graph maps.
The graph $U=A$ is defined there over $e'=1$, and is identified with
$H_{48}|_{e'=1}$. Its comparison is the restriction of (A.16.2)
over $b=0$. Thus all three \S51 groups are related to the earlier
groups by these explicit base changes, character maps, and central
quotients. The resulting schemes are not isomorphic: their generic
relative dimensions are respectively two and three, and their bad
fibres have the distinct descriptions proved below.
\end{statement}
\begin{proof}
Solving $A\rho=\rho A$ and $B\sigma=e\sigma B$ gives the two
displayed matrix forms, and their determinants give exactly the
stated equations. The character map is fppf surjective: given its
three coordinates, adjoining square roots of the required unit norms
and taking scalar elements in the quadratic algebras supplies lifts.
The extensions adjoining square roots are finite free and faithfully
flat, including in characteristic two. Its kernel is precisely the
two norm-one groups. Formula (A.16.1) is a homomorphism by the ring
and sign formulas in A.5--A.7. Over $\mathbb Z[1/6]$, its inverse
recovers the constant signs and sets $k=(\mu-r_0)/12$, proving the
asserted isomorphism.

The base-changed character map retains its exact kernel, so its
composition with the group section of $\Gamma/\mu_\rho$ has the
kernel and image claimed in (A.16.2). For the full matrix triple,
the additional determinant-prescribed $U$ coordinate has kernel
$\mathrm{SL}_2$. The section $U=\operatorname{diag}(\det B,1)$
gives the graph at zero, and $U=B$ gives the graph at one-half.
For $U=A$, the determinant condition reads $e'=1$, proving exactly
the specified restriction. This derives the comparison from the
equations used in LM1 B.2--B.5 instead of importing their conclusion.

Here is also a direct fibre computation. Eliminating the common
determinant gives
\[
 H_{48}=\operatorname{Spec}
 \mathbb Z[c,d,t,u,e,e',Q_2^{-1}]/
 (Q_3-ee'Q_2,e^2-1,(e')^2-1).
\]
Relative to the sign scheme $\mu_2^2$ its hypersurface is smooth:
the gradient in $c,d,t,u$ is
$(2c+d,c+2d,-2ee't,-2ee'u)$ and cannot vanish on $Q_2\ne0$.
For characteristic different from two, use the last two entries;
in characteristic two the first two become $d,c$ and cannot both
vanish because $Q_3$ is a unit. The geometric fibres are nonempty,
so this is a smooth surjective map of relative dimension three.
Thus the generic dimension is three. At three,
$Q_3=(c-d)^2$; setting $v=c-d\ne0$ shows four smooth geometric
components $\mathbb G_m^2\times\mathbb G_a$.
At two put $x=e-1,y=e'-1$. The sign algebra is
$\mathbb F_2[x,y]/(x^2,y^2)$; after reduction, $Q_2=(t+u)^2$ and
the remaining equation is the norm-one equation
$a^2+ab+b^2=1$. Hence the reduced fibre is
\[
 \mathbb G_a\times\mathbb G_m\times
 R^1_{\mathbb F_4/\mathbb F_2}\mathbb G_m.
\]
It is geometrically connected, and smoothness over the sign algebra
gives generic length four before reduction. Restricting $e'=1$
replaces that length by two. The full $\mathcal M_{48}$ is, as a
scheme, the product of this scheme with $\mathrm{SL}_2$, by
$U=\operatorname{diag}(\det B,1)h$.
These fibre calculations and the explicit morphisms above prove
the asserted comparison with A.15, whose characteristic-two fibre
has eight reduced affine-plane components. They also prove flatness
of the \S48 models and their failure of smoothness exactly at two.
No nonexistence of further mathematical comparisons is asserted.
\end{proof}

\begin{statement}{A.17}{Proved}{Linear congruence descent from the actual braid automorphism}
Use the standard tangential and braid normalization
\[
 B_3=\langle\beta_1,\beta_2:
              \beta_1\beta_2\beta_1=\beta_2\beta_1\beta_2\rangle,
 \quad x=\beta_1^2,\quad y=\beta_2^2.
\]
For $(\mu,f)\in\widehat{\mathrm{GT}}_0$, where
$f\in\widehat{\langle x,y\rangle}'$, let
\[
 \Phi(\beta_1)=\beta_1^\mu,\qquad
 \Phi(\beta_2)=f^{-1}\beta_2^\mu f.                \tag{A.17.1}
\]
For every $n\geq2$, define the continuous surjection
\[
 q_n:\widehat B_3\longrightarrow\mathrm{SL}_2(\mathbb Z/n),
 \quad\beta_1\mapsto U=\begin{pmatrix}1&1\\0&1\end{pmatrix},
 \quad\beta_2\mapsto W=\begin{pmatrix}1&0\\-1&1\end{pmatrix}.
\]
There is a unique matrix
\[
 V_n(\Phi)=\begin{pmatrix}\mu&t_n(\Phi)\\0&1\end{pmatrix}
 \quad\text{such that}\quad
 q_n\Phi=\operatorname{Int}(V_n(\Phi))q_n.          \tag{A.17.2}
\]
The maps $\Phi\mapsto V_n(\Phi)$ are continuous homomorphisms,
compatible under reduction. In particular all linear congruence
kernels are invariant. They give a specified action on the full
$\mathrm{GL}_2(\mathbb Z/n)$ by the same conjugation formula.
This supplies a direct proof of the descent needed in LM2 C.d1 and
its requested independent check at $n=12$.
\end{statement}
\begin{proof}
The group-theoretic input is that (A.17.1) is an automorphism of
$\widehat B_3$ for an actual $\widehat{\mathrm{GT}}_0$ pair.
Lochak--Schneps prove the braid identity in Lemma 5 of
\cite{A_LochakSchneps}; Harbater--Schneps, \S1.3, Proposition 4,
proof of (i)$\Rightarrow$(iii), explicitly applies that lemma to the
full profinite group, \cite{A_HarbaterSchneps}. Both proofs were
opened and checked. The input uses the two cycle relations and the
automorphism requirement; no pentagon is used here.
The following reduction from that theorem to (A.17.2) is elementary
and also proves the requisite generation condition cannot be omitted.

Write $A=q_n\Phi(\beta_1)=U^\mu$ and
$B=q_n\Phi(\beta_2)=\left(\begin{smallmatrix}a&b\\c&d\end{smallmatrix}\right)$.
Conjugacy to $W^\mu$ gives $a+d=2$, $ad-bc=1$.
The lower-left entry of $ABA=BAB$ gives
\[
 c(1+\mu c)=0.                                    \tag{A.17.3}
\]
The pair $A,B$ generates all of $\mathrm{SL}_2(\mathbb Z/n)$,
because $q_n$ is onto and $\Phi$ is an automorphism. For each prime
$p\mid n$, the element $c\bmod p$ cannot vanish: otherwise both
matrices would be upper triangular, and could not generate
$\mathrm{SL}_2(\mathbb F_p)$. Thus $c$ is a unit modulo $n$.
Equation (A.17.3) yields $c=-\mu^{-1}$; the determinant and trace
then give
\[
 d=2-a,\quad b=\mu(a-1)^2,\quad t_n=\mu(1-a).
\]
Direct multiplication proves
$V_nUV_n^{-1}=A$, $V_nWV_n^{-1}=B$. Agreement on the topological
generators proves (A.17.2) on the full profinite group.
For surjectivity of $q_n$, elementary row operations over each local
ring $\mathbb Z/p^r$ reduce a determinant-one matrix to diagonal
form, and diagonal determinant-one matrices are products of
elementary matrices. Chinese remainders combine these operations.
Powers of $U,W$ give all the elementary matrices.

A matrix commuting with $U,W$ is scalar, by comparing its four
entries. A scalar relating two matrices with lower-right entry one
is the identity. This proves uniqueness. The product $V_n(\Phi)
V_n(\Psi)$ has the same normalization and implements the composite
automorphism, which proves the homomorphism assertion. The formula
for $a$ in terms of $q_n(f)$ proves continuity; uniqueness proves
compatibility in $n$. Taking the same inner conjugation on
$\mathrm{GL}_2(\mathbb Z/n)$ is now a well-typed extension of the
action on its determinant-one subgroup.

For comparison with the modular group, the central braid
$z=(\beta_1\beta_2)^3$ maps to $\omega$, and
$\widehat B_3/\overline{\langle z^2\rangle}
\simeq\widehat{\mathrm{SL}_2(\mathbb Z)}$.
The cited proof also gives $\Phi(z)=z^\mu$, so this quotient is
preserved. The subgroup $\widehat\Pi_0$ is the inverse image of
$V_4$ modulo four, preserved by the matrices in A.18.
Alternatively its generators are
$\lambda_0=\omega U^{-2}$, $\lambda_1=\omega W^{-2}$; their images
are respectively $\lambda_0^\mu$ and a conjugate of
$\lambda_1^\mu$ by an element of $\widehat\Pi_0'$. To verify the
last assertion, evaluate the commutator word $f$ in
$U^2=\omega\lambda_0^{-1}$ and
$W^2=\omega\lambda_1^{-1}$: its central factors cancel. This
constructs the comparison to the action on $M/12M$ used below.

The program \texttt{a\_finite.py} independently enumerates every
signed matrix-coboundary solution used in LM2 C.d1 at level twelve:
there are twelve for each of the four multipliers and no failed
intertwiner. The program \texttt{a\_braid.py} additionally
enumerates the eight possible values of $f$ in the finite
commutator image, for every multiplier. All thirty-two pairs
satisfy the braid relation; eight pairs fail generation and are
explicitly rejected. The remaining twenty-four pairs give the
twelve distinct matrices in A.18. This records both the computation
and the exact reason the discarded matrix solutions cannot be
automorphisms of the profinite source.
\end{proof}

\begin{statement}{A.18}{Proved}{Exact images at level twelve, with the genuine normalization stated}
For the normalization (A.17.1)--(A.17.2), the images of
$G_{\mathbb Q}$, $\widehat{\mathrm{GT}}$, and
$\widehat{\mathrm{GT}}_0$ in $\mathrm{GL}_2(\mathbb Z/12)$ are all
\[
 \mathcal V_{12}=
 \left\{\begin{pmatrix}\mu&t\\0&1\end{pmatrix}:
     \mu\in\{1,5,7,11\},\quad t\in\{0,4,8\}\right\},
 \qquad |\mathcal V_{12}|=12.                      \tag{A.18.1}
\]
Their specified actions in $\operatorname{Aut}(\mathrm{GL}_2(\mathbb Z/12))$
therefore have image $\operatorname{Ad}(\mathcal V_{12})$ of order
$12$. Their images in
$\operatorname{Out}(\mathrm{GL}_2(\mathbb Z/12))$ are trivial.
The canonical outer images on $\mathrm{SL}_2(\mathbb Z/12)$ are all
$(\mathbb Z/12)^*/((\mathbb Z/12)^*)^2=C_2^2$, of order four.
On $M/12M$ the genuine image of each group is the scalar subgroup
$\{\mu I:\mu\in(\mathbb Z/12)^*\}$, also of order four.

For $G_{\mathbb Q}$, choose the tangential base point at $\lambda=0$
with tangent $+1$ and the real interval path to the tangent at one;
this is the Belyi normalization in (A.17.1).
The kernel of its normalized level-twelve matrix representation is
exactly
\[
 G_{\mathbb Q(\zeta_{12},\sqrt[3]{2})}.           \tag{A.18.2}
\]
The extended group generated by the geometric modular group and
any of the three arithmetic groups acts on the same linear
congruence quotient with image
$\operatorname{Ad}(\mathrm{GL}_2(\mathbb Z/12))$ of order $1152$.
These statements give the precise actions requested after LM2;
the order $32$ remains the normal subgroup $K_{32}$ of A.3.
\end{statement}
\begin{proof}
Modulo four, $q_4(x)=U^2$ and $q_4(y)=W^2$ commute, so
$q_4(f)=I$. Hence $q_4\Phi(\beta_2)=W^\mu$, whose upper-left
entry is one. The formula in A.17 gives $t\equiv0\pmod4$.
Thus the whole $\widehat{\mathrm{GT}}_0$ image lies in
$\mathcal V_{12}$. This set is a group under matrix multiplication
and has twelve elements. To prove equality for the arithmetic
subgroup, we now compute a compatible finite geometric cover,
including its base-point comparison.

Consider the Legendre curve
$E_\lambda:y^2=x(x-1)(x-\lambda)$ over
$\mathbb P^1-\{0,1,\infty\}$, and its finite etale cover of
cyclic order-three subgroups. Its four points are the four
$x$-coordinates of nonzero three-torsion modulo sign. With the
oriented vanishing cycles above the intervals $[0,\lambda]$ and
$[\lambda,1]$, a positive loop around either colliding pair of
branch points is a full twist. On the double cover it acts by the
square of the transvection in the corresponding vanishing cycle.
Taking their intersection to be one gives the matrices
$U^2$ and $W^2$. This constructs the geometric monodromy comparison
from $x,y$ in (A.17.1) to the projectivization of $E[3]$.
Changing the simultaneous homology basis conjugates both sides;
the four-point permutation representation and its coefficient
field are unchanged. The central signs used to define $\lambda_i$
disappear under this specified projectivization.

We also specify the arithmetic fibre functor. At the tangent
$\lambda=0,+1$, the fibre is obtained by adjoining formal roots
$\lambda^{1/m}$ and letting $G_{\mathbb Q}$ act on their coefficients,
fixing the symbols $\lambda^{1/m}$. The real interval path gives the
usual second generator. In every abelian Kummer covering generated
by $\lambda^{1/m}$ and $(1-\lambda)^{1/m}$, the leading coefficients
at the two tangents are one. Thus the path cocycle has zero
abelianization; the induced based Galois action has precisely
$x\mapsto x^\chi$, $y\mapsto f^{-1}y^\chi f$ with $f\in\widehat F_2'$.
The arithmetic compatibility and the inclusions
$G_{\mathbb Q}\subset\widehat{\mathrm{GT}}\subset
\widehat{\mathrm{GT}}_0$ are the tangential actions described in
\cite[\S0.3 and \S1.2--1.3]{A_HarbaterSchneps}.
Consequently the arithmetic conjugation of this four-point
Legendre monodromy is the projective reduction of $V_{12}$.
There is no extra projective conjugator ambiguity here: a matrix
centralizing the projective images of $U^2,W^2$ modulo three
centralizes their actual matrices, since their trace two is
nonzero; comparing entries then makes it scalar.

The three-division polynomial, obtained from the tangent-and-chord
addition formula on the displayed cubic, is
\[
 \psi_3(x)=3x^4-4(1+\lambda)x^3+6\lambda x^2-\lambda^2. \tag{A.18.3}
\]
At zero it is $x^3(3x-4)$. The root with constant term $4/3$
lies in $\mathbb Q[[\lambda]]$, by the formal implicit-function
theorem. For the other roots put $\lambda=s^3$, $x=s^2z$.
Dividing the polynomial by $s^6$ gives
\[
 3s^2z^4-4(1+s^3)z^3+6sz^2-1.                    \tag{A.18.4}
\]
Each root $c$ of $c^3=-1/4$ is simple at $s=0$, and uniquely
lifts to $z\in\mathbb Q(c)[[s]]$. Hence the exact field of all
Puiseux coefficients is
$L=\mathbb Q(\zeta_3,\sqrt[3]{2})$: its containment follows by
the unique recursion for the coefficients, and the reverse
containment follows from the three leading coefficients, because
$( -2c)^3=2$. The polynomial $X^3-2$ is irreducible by Eisenstein
at two, and its discriminant $-108$ is not a square. Therefore
$[L:\mathbb Q]=6$ and the arithmetic action on the four lines
has image $S_3$, fixing the one rational branch.

Modulo four, $V_{12}$ is $\operatorname{diag}(\chi,1)$, so its
remaining field is $\mathbb Q(i)$. The field $L$ has unique
quadratic subfield $\mathbb Q(\sqrt{-3})$, by the subgroup lattice
of $S_3$, and therefore $L\cap\mathbb Q(i)=\mathbb Q$.
The pair consisting of $V\bmod4$ and $[V\bmod3]$ determines
every $V\in\mathcal V_{12}$: a scalar with lower-right entry one
is the identity, and Chinese remainders recover both entries.
Thus the arithmetic kernel fixes exactly the degree-twelve field
$L(i)=\mathbb Q(\zeta_{12},\sqrt[3]{2})$, proving (A.18.1)--(A.18.2).
The two containing GT groups have the same image by the already
proved upper bound. The formal polynomial calculations (A.18.3)--
(A.18.4), the leading field, and all twelve matrices are checked
in \texttt{a\_braid.py}.

The centralizer of $\mathrm{SL}_2(\mathbb Z/12)$ in
$\mathrm{GL}_2(\mathbb Z/12)$ consists of its four scalar units.
The normalization $V_{22}=1$ makes their intersection with
$\mathcal V_{12}$ trivial. This proves the genuine automorphism
image count and the trivial outer image on the full GL group.
A conjugation by $V$ is inner on SL exactly when
$V\in\mathrm{SL}_2(\mathbb Z/12)\cdot
((\mathbb Z/12)^*I)$; taking determinants gives its unit square
class. All four units modulo twelve square to one, proving the
outer SL count. The scalar action on $M/12M$ follows from the
generator comparison at the end of A.17, and cyclotomic
surjectivity realizes all four scalars. Finally the geometric
image is SL and the determinants of $\mathcal V_{12}$ are all
units, so $\langle\mathrm{SL}_2(\mathbb Z/12),\mathcal V_{12}\rangle
=\mathrm{GL}_2(\mathbb Z/12)$. Quotienting its order $4608$ by
the scalar centralizer of order four gives $1152$.

This specifies which choices give a genuine action and which
quotients remove them. The congruence action is obtained from
$\widehat B_3$ and its modular quotient; A.4 proves that the
reduction to the unquotiented GL or SL group cannot instead be
defined by first passing through $\mathfrak S_{\rm d}$.
No equality between the normalized upper-right entry and the
Kummer cocycle based at $1/2$ is used in this proof.
\end{proof}

\begin{statement}{A.19}{Proved}{The normalized cyclotomic arithmetic image in the minimum approximation}
In the tangential normalization of A.17--A.18, the arithmetic
homomorphism into the final \S51 minimum approximation has image in
$\mathcal M(0)$: its effective coordinate is $U=u_{\underline\mu}$.
Its restriction to $G_{\mathbb Q^{\rm cyc}}$ has image in the abelian
group
\[
 \ker(\mathcal M(0)(\widehat{\mathbb Z})^{\Delta}
                   \longrightarrow\widehat{\mathbb Z}^{*})
       =C_2\times\widehat{\mathbb Z}.              \tag{A.19.1}
\]
After the central quotient by $\mu_\rho$, this kernel is
$\widehat{\mathbb Z}$ under addition. In particular the closed
commutator subgroup of $G_{\mathbb Q^{\rm cyc}}$ maps to the identity
in this normalized minimum approximation. Here the superscript
$\Delta$ requires the single global component choice described in
A.1; it is essential to (A.19.1).
\end{statement}
\begin{proof}
The congruence-normalized braid action has conjugator
$f\in\widehat F_2'$. In the modular quotient its image is in
$\widehat\Pi_0'$ by the evaluation calculation in A.17, so it
vanishes in the minimum approximation. Thus both elementary
generators $U,W$ are sent to their $\mu$-th powers.
The automorphism (A.8.5) has exactly these images on diagonal
profinite points: from $U=(\sigma\rho)^{-1}$,
$W=(\rho\sigma)^{-1}$ and
$\lambda_0=\omega U^{-2}$, $\lambda_1=\omega W^{-2}$,
the possible central factor is $\omega^{\nu+a}=1$, since
$\nu=(\mu-1)/2\equiv a\pmod2$.
This proves $U=u_{\underline\mu}$ for the effective coordinate,
which is precisely the graph defining $\mathcal M(0)$.
The retained normalizer conjugators give its other two coordinates
as in (183), so its projection to those coordinates is the
isomorphism of A.15.

For a diagonal cyclotomic point with $\mu=1$, its residue is one
and the unique quotient $(\mu-1)/12$ is zero in
$\widehat{\mathbb Z}$: hence $a=b=k=0$. Formula (A.14.2) on this
fibre is $(i,s)(i',s')=(i+i',s+s')$. This proves (A.19.1), its
central quotient, and the commutator assertion. It also explains
why replacing diagonal profinite points by arbitrary
$\gamma(\mathbb F_2)$-points would give the wrong kernel.
\end{proof}
